% Shared section: data on national well-being, income and regions (WVS, Gallup, WHR, WDI)
\subsection{World Values Survey}\label{subsec:wvs}

The World Values Survey (WVS) is the longest-running cross-national survey of values and subjective well-being. I use the WVS time-series file \citep{wvs_trend_2022}, which pools the seven waves conducted between 1981 and 2022: \nWVSRespondents respondents in \nWVSCountryYears country-years and \nWVSCountries countries. National samples are designed to be representative of the adult population, typically with 1,000 to 3,000 respondents, and most interviews are conducted face-to-face (though some recent national surveys used web, mail or mixed modes, e.g. the United States in 2017 and Japan in 2019). All statistics use the survey weights provided by the WVS.

The WVS contains two questions on subjective well-being:
\begin{itemize}
    \item \textit{Happiness}: ``Taking all things together, would you say you are: \textit{Very happy}; \textit{Quite happy}; \textit{Not very happy}; \textit{Not at all happy}.''
    \item \textit{Life satisfaction}: ``All things considered, how satisfied are you with your life as a whole these days?'', answered on a scale from 1 (\textit{Completely dissatisfied}) to 10 (\textit{Completely satisfied}).
\end{itemize}
Non-responses are rare (\nonresponseSatisfactionWVS\% for satisfaction and \nonresponseHappinessWVS\% for happiness on average) and uncorrelated with GDP per capita (correlation: \corNonresponseGDP); they are excluded.

\ifshowregion
\paragraph{Indicators of national well-being.}
Because there is no consensus on how to aggregate ordinal answers into a national indicator, and because the ranking of groups can depend on this choice \citep{bond_sad_2019}, I compute eight indicators for each country-year: \textit{Happy} (share answering \textit{Quite} or \textit{Very happy}); \textit{Very Happy}; \textit{Very Unhappy} (share \textit{Not at all happy}); \textit{V.~Happy -- V.~Unhappy} (difference of the two previous shares); \textit{Happiness (mean)} (answers coded $+3$, $+1$, $-1$, $-3$); \textit{Satisfaction (mean)}; \textit{Satisfied} (share answering 6 to 10); and \textit{Happy + Satisfied} (average of \textit{Happy} and \textit{Satisfied}), the index used by \citet{inglehart_genes_2000}. I also compute the share of people whose satisfaction is below 60\% of the national mean, by analogy with relative poverty lines. The indicators are positively but imperfectly correlated: across country-years, the correlation between \textit{Happy} and \textit{Satisfied} is \corHappySatisfied, and between \textit{Very Happy} and \textit{Satisfaction (mean)} it is only \corVeryHappySatisfiedMean.
\fi

\subsection{Gallup World Poll and World Happiness Report}\label{subsec:gallup}

The Gallup World Poll interviews around 1,000 adults per country and year in more than 140 countries, by telephone where coverage is high (typically in high-income countries) and face-to-face elsewhere \citep{gallup_methodology_2022}. Its life-evaluation question is the Cantril ladder \citep{cantril_pattern_1965}: ``Please imagine a ladder, with steps numbered from 0 at the bottom to 10 at the top. The top of the ladder represents the best possible life for you and the bottom of the ladder represents the worst possible life for you. On which step of the ladder would you say you personally feel you stand at this time?''

I use the distribution of answers by country and wave, tabulated from the Gallup World Poll microdata (\nGallupCountryYears country-years, \nGallupCountries countries, 2006--2023), from which I compute the same satisfaction indicators as for the WVS (mean, shares answering 6--10, 8--10, 9--10, 10, 0--4 and 0--2). These distributions are unweighted. Gallup's wave $w$ corresponds to year $w+2005$: with this mapping, three-year averages of the ladder computed from these distributions correlate at \corGallupWHR with the (weighted) three-year averages published in the World Happiness Report, against \corGallupWHRoldMapping with the mapping used in earlier versions of this project (year $w+2000$). To cover the most recent years, I also use the ladder means published in the World Happiness Report 2026 \citep{whr_2026}, which average Gallup data over three years up to 2023--2025.

\subsection{Income}\label{subsec:income}

My benchmark income measure is the log of GDP per capita at purchasing power parity (PPP, constant 2021 international dollars) from the World Development Indicators \citep{worldbank_wdi_2026}. As an alternative, I use GDP per capita at market exchange rates (constant 2015 US dollars). Missing values are filled in automatically, in three steps: (i) PPP series, which start in 1990, are extended backwards (and gaps filled) using the growth rate of real GDP per capita; (ii) for countries absent from the World Bank data (e.g. Taiwan, or Venezuela in recent years), I use the IMF World Economic Outlook \citep{imf_weo_2026}, converted into constant dollars using the U.S. ratio between the World Bank and IMF series; (iii) remaining gaps are filled with the closest year available (Montenegro 1996). Territories without national accounts use the GDP of the closest entity (e.g. the United Kingdom for Northern Ireland). A robustness check excludes all imputed values.

\ifshowregion
Following the previous version of this analysis, I also use discrete income indicators, computed within each estimation sample: income sextiles and income clusters obtained by $k$-means clustering of log GDP per capita, with $k = 5$, 6 or 7. The clusters let the data choose the thresholds between income groups and allow for non-linear relationships, at the cost of more parameters than the log-linear specification.
\fi

\ifshowregion
\subsection{Regions}\label{subsec:regions}

My main classification groups countries into five world regions inspired by the UN regional groups: \textit{Africa}, \textit{Asia} (including the Middle East and Central Asia), \textit{Eastern Europe} (the former Eastern Bloc, excluding Central Asia), \textit{Latin America} (including the Caribbean), and \textit{Western} countries (Western Europe, North America, Australia, New Zealand, and Turkey). For robustness, I use four alternative classifications: a six-region variant separating the Middle East and grouping Central Asia with the former Eastern Bloc; the seven World Bank regions; the five continents; and the UN geoscheme sub-regions (about twenty groups).
\fi
